% Part: methods
% Chapter: proofs
% Section: inference-patterns

\documentclass[../../../include/open-logic-section]{subfiles}

\begin{document}

\olfileid{mth}{prf}{pat}

\olsection{د استنتاج بڼې}

ثبوتونه له بېلو بېلو استنتاجونو جوړېږي۔ کله چې استنتاج
کوو، عموماً يې په «نو»، «له دې امله»، يا «ځکه نو» غوندې
ټکي څرګندوو۔ استنتاج ډېر ځله په هغو يو يا دوو خبرو
ولاړ وي چې په ثبوت کښې مو لا په لاس کښې دي: ښايي
فرض مو کړې وي، يا مو مخکې له بل استنتاج څخه
اخيستې وي۔ د روڼوالي دپاره دغو خبرو ته نښې
ورکولے شو، او د استنتاج پر مهال وايو چې کومې
نورې خبرې کاروو۔ استنتاج ډېر ځله دا هم تشريح کوي
چې نوې پايله مو له مخکېنيو خبرو څخه \emph{ولې}
راوځي۔ په ثبوتونو کښې د استنتاج ځينې عامې بڼې
بيا بيا کارېږي؛ لاندې يې څو بڼې څېړو۔ ځينې
بڼې، لکه د استقرا له لارې ثبوت، زياتې پېچلې دي
او وروسته به ورباندې خبرې وکړو۔

د استنتاج يوه بڼه مو لا څېړلې ده: د تعريف پرانيستل
يا کارول۔ د تعريف په پرانيستلو کښې هغه خبره چې
تعريفېدونکے پکې راغلے وي، د تعريفوونکي بيان په
کارولو بيا وايو۔ د بېلګې په توګه، فرض کړئ چې
د ثبوت په ترڅ کښې مو $D = E$ لا ثابته کړې ده
(الف)۔ بيا د سټونو دپاره د $=$ تعريف کاروو
او استنتاج کوو: «نو، له (الف) څخه د تعريف له مخې،
د~$D$ هر !!{element} د~$E$ !!a{element} دے، او برعکس»۔

يو څه مغشوشوونکې دا ده چې موږ د استنتاج توجيه
ډېر ځله د هماغه استنتاج د کولو پر مهال نه،
بلکې تر هغه مخکې ليکو۔ فرض کړئ $D = E$ مو لا
نه دي ثابت کړي، خو ثابتول يې غواړو۔ که $D = E$
زموږ غوښتل شوې پايله وي، هدف د تعريف په کارولو
بيا ويلے شو: د $D = E$ د ثابتولو دپاره بايد
ثابته کړو چې د~$D$ هر !!{element} د~$E$
!!a{element} دے، او برعکس۔ نو د ثبوت بڼه به
داسې وي: (الف) د~$D$ هر !!{element} د~$E$
!!a{element} ثابت کړئ؛ (ب) د~$E$ هر
!!{element} د~$D$ !!a{element} ثابت کړئ؛
(ج) ځکه نو، له (الف) او (ب) څخه د $=$ د تعريف
له مخې، $D = E$۔ خو عموماً يې داسې نه ليکو۔
پر ځاے يې ښايي داسې وليکو:
\begin{quote}
غواړو وښيو چې $D = E$۔ د~$=$ له تعريفه، د دې
دپاره بايد وښيو چې د~$D$ هر !!{element} د~$E$
!!a{element} دے، او برعکس۔

(الف) \dots (دا ثبوت چې د~$D$ هر !!{element}
د~$E$ !!a{element} دے) \dots

(ب) \dots (دا ثبوت چې د~$E$ هر !!{element}
د~$D$ !!a{element} دے) \dots
\end{quote}

\subsection{د عطف کارول}

ښايي د استنتاج تر ټولو ساده بڼه دا وي چې د عطف
له دوو غړو څخه يو د پايلې په توګه واخلو۔ يعنې که
مو فرض کړي يا مخکې ثابت کړي وي چې $p$ او~$q$،
نو $p$ (او همدارنګه~$q$) استنتاجولے شو۔ دا
دومره بنسټيز استنتاج دے چې ډېر ځله يې جلا يادونه
نه کېږي۔ د بېلګې په توګه، کله چې د $D = E$ تعريف
پرانيزو، دا مو ثابته کړې وي چې د~$D$ هر
!!{element} د~$E$ !!a{element} دے، او برعکس۔
له دې څخه دا اخيستلے شو چې د~$E$ هر
!!{element} د~$D$ !!a{element} دے؛ همدا د
«برعکس» برخه ده۔

\subsection{د عطف ثابتول}

کله کله هغه خبره چې ثابتول يې غواړئ د عطف
بڼه لري: له تاسو څخه غوښتل کېږي چې «$p$ او
$q$ ثابت کړئ»۔ دلته يوازې دوه کارونه کوئ:
لومړے $p$ ثابت کړئ، بيا $q$ ثابت کړئ۔ ثبوت
په دوو برخو وېشلے شئ او د روڼوالي دپاره يې
ونوموئ۔ د لومړنيو يادښتونو پر مهال ښايي د پاڼې
په سر کښې «(۱) $p$ ثابت کړئ» او د پاڼې په منځ
کښې «(۲) $q$ ثابت کړئ» وليکئ۔ (کېدے شي له
تاسو څخه د عطف د ثابتولو غوښتنه په څرګنده
نه وي شوې، خو د ثبوت دپاره يې ثابتول اړين
ومومئ۔ د بېلګې په توګه، که د $D = E$ د ثابتولو
غوښتنه وي، د~$=$ د تعريف له پرانيستلو وروسته
بايد دا دواړه ثابت کړئ: د~$D$ هر !!{element}
د~$E$ !!a{element} دے \emph{او} د~$E$ هر
!!{element} د~$D$ !!a{element} دے۔)

\subsection{د فصل ثابتول}

که ثابتېدونکې خبره د فصل بڼه ولري (يعنې
«$p$ يا $q$» وي)، د دوو بديلو غړو له ډلې
د يوه رښتياوالے ښودل بس دے۔ خو په عمل کښې
تقريباً هېڅکله داسې نه وي چې يو بديل غړے
په نېغه د قضيې له فرضونو راووځي۔ ډېر ځله
د قضيې فرضونه پخپله په فصل کښې وي، يا تاسو
ښيئ چې د يوه ډول هر څيز له دوو ځانګړنو څخه
يوه لري، خو ځينې څيزونه لومړۍ او نور دويمه
ځانګړنه لري۔ همدلته د حالتونو له مخې ثبوت
ګټور وي (لاندې يې وګورئ)۔

\subsection{شرطي ثبوت}

ډېرې قضيې چې ورسره مخامخ کېږئ شرطي بڼه
لري (يعنې وښيئ چې که $p$ وي، نو $q$ هم
رښتيا دے)۔ د دغه ډول ثبوت برابرول اسان دي:
د شرطي مقدم فرض کړئ (دلته $p$)، او له هغه
څخه پايله~$q$ ثابته کړئ۔ نو که قضيه وايي
«که $p$، نو $q$»، ثبوت په «$p$ فرض کړئ»
پيلوئ او په پای کښې بايد~$q$ مو ثابت کړے وي۔

شرطي خبرې په بېلابېلو لارو ويل کېږي۔ نو د
«که $p$، نو $q$» پر ځاے قضيه ښايي ووايي:
«$p$ يوازې که $q$»، «$q$ که $p$»، يا
«$q$، په دې شرط چې $p$»۔ دا ټول يوه معنا
لري: $p$ فرض کړئ او له هغه څخه~$q$ ثابت
کړئ۔ په ياد ولرئ چې دوه اړخيز شرطي («$p$
که او يوازې که $q$») په حقيقت کښې دوه
شرطي خبرې دي: که $p$، نو $q$؛ او که $q$،
نو~$p$۔ نو دوه شرطي ثبوتونه پکار دي: يو
د لومړي لوري او بل د دويم لوري دپاره۔
کله کله بيا د «که او يوازې که» خبره د
همداسې دوه اړخيزو خبرو په لړۍ ثابتولے شئ،
چې له «$p$» څخه پيل او «$q$» ته رسېږي؛
خو هر ګام بايد په رښتيا د دواړو لورو
شرط وي۔

\subsection{عمومي دعوې}

د عمومي دعوې کارول اسان دي: که خبره د
هر څيز دپاره رښتيا وي، د هر ځانګړي څيز
دپاره هم رښتيا ده۔ نو، د بېلګې په توګه،
که ستاسو د ثبوت فرض $A \subseteq B$ وي،
د~$\subseteq$ د تعريف له پرانيستلو سره
دا معنا لري چې د هر $x \in A$ دپاره
$x \in B$۔ ځکه نو که لا پوهېږئ چې
$z \in A$، دا پايله اخيستلے شئ چې
$z \in B$۔

د عمومي دعوې ثابتول ښايي لږ ستونزمن
وبرېښي۔ دا ډول خبرې عموماً داسې وي:
«که $x$ د~$P$ ځانګړنه ولري، نو د~$Q$
ځانګړنه هم لري»، يا «ټول $P$s د $Q$s
دي»۔ البته تل به يې بڼه کټ مټ داسې نه
وي؛ د پوښتنې د کره غوښتنې پېژندل لږ
تمرین غواړي۔ خو ډېر ځله بايد ثابته کړو
چې د يوې ځانګړنې لرونکي ټول څيزونه
يوه بله ځانګړنه هم لري۔

د عمومي دعوې د ثابتولو لار دا ده چې
د لومړۍ ځانګړنې لرونکو څيزونو دپاره
نومونه يا متغيرونه راوړئ، او بيا وښيئ
چې دويمه ځانګړنه هم لري۔ ويلے شو چې
د \emph{ټولو}~$P$s دپاره د خبرې د
ثابتولو دپاره بايد يې د يوه
\emph{هر ډول}~$P$ دپاره ثابت کړئ۔
راوړل شوے نوم د همداسې هر ډول~$P$
نوم دے۔ عموماً د دغو هر ډول څيزونو
د نومونو دپاره يو توری کاروو، او
توري له دود سره سم ټاکو: لکه $n$
د طبيعي عددونو، $!A$ د !!{formula}s،
$A$ د سټونو، او $f$ د تابعو دپاره۔

اصل چل دا دے چې په ټول ثبوت کښې
عموميت وساتئ۔ په پيل کښې فرض کړئ چې
يو هر ډول څيز («$x$») د~$P$
ځانګړنه لري؛ بيا يوازې د تعريفونو
او اجازه شوو فرضونو له مخې وښيئ
چې $x$ د~$Q$ ځانګړنه هم لري۔
ځکه مو د $x$ په اړه، له $P$ لرلو
پرته، بل ځانګړے شرط نه دے ايښے،
ويلے شئ چې هر هغه څيز چې $P$
لري، $Q$ هم لري۔ لنډه دا چې $x$
د~$P$ ځانګړنې د \emph{ټولو}
څيزونو ځايناستے دے۔

\begin{prop}
  د هر دوو سټونو $A$ او $B$ دپاره، $A \subseteq A \cup B$۔
\end{prop}

\begin{proof}
  $A$ او $B$ دې هر ډول سټونه وي۔ غواړو وښيو چې $A
  \subseteq A \cup B$۔ د $\subseteq$ له تعريفه دا
  معنا لري: د هر $x$ دپاره، که $x \in A$ وي، نو
  $x \in A \cup B$۔ نو $x \in A$ دې د~$A$ يو
  هر ډول !!{element} وي۔ بايد وښيو چې $x \in A
  \cup B$۔ څرنګه چې $x \in A$، نو $x \in A$
  يا $x \in B$۔ له دې څخه $x \in
  \Setabs{x}{x \in A \lor x \in B}$۔ خو دا،
  د $\cup
  $ له تعريفه، $x \in A \cup B$ معنا لري۔
\end{proof}

\subsection{د حالتونو له مخې ثبوت}

فرض کړئ فصل د يوه فرض يا مخکې ثابتې شوې پايلې
په توګه لرئ: فرض مو کړي يا ثابت کړي دي چې $p$
يا $q$ رښتيا دے۔ غواړئ $r$ ثابت کړئ۔ دا کار
په دوو پړاوونو کښې کوئ: لومړے $p$ رښتيا
فرض کړئ او~$r$ ثابت کړئ؛ بيا $q$ رښتيا
فرض کړئ او~$r$ يو ځل بيا ثابت کړئ۔ دا
لار ځکه کار کوي چې له دوو بديلو څخه
يو رښتيا فرض شوے يا پېژندل شوے دے۔
دواړه پړاوونه ښيي چې هر يو په جلا
توګه د~$r$ د رښتياوالي دپاره بس دے۔
(که دواړه رښتيا وي، د $r$ د رښتياوالي
دپاره يو نه، بلکې دوه دليلونه لرو۔
اړينه نه ده چې $r$ د $p$ او~$q$
دواړو په يو وخت فرضولو سره بيا جلا
ثابت کړو۔) د خپلې لارې د څرګندولو
دپاره وايو چې «حالتونه بېلوو»۔
د بېلګې په توګه، فرض کړئ پوهېږو
چې $x \in B \cup C$۔ $B \cup C$
د $\Setabs{x}{x \in B \text{ or } x \in C}$
په توګه تعريف شوے دے۔ يعنې د تعريف
له مخې، $x \in B$ يا $x \in C$۔
له دې څخه $x \in A$ داسې ثابتوو:
لومړے $x \in B$ فرض کړئ او له هغه
څخه $x \in A$ ثابت کړئ؛ بيا
$x \in C$ فرض کړئ او له هغه
څخه يو ځل بيا $x \in A$ ثابت
کړئ۔ د فرض تر لاندې «حالتونه
بېلوو» وليکئ، بيا «حالت (۱):
$x \in B$»؛ او د پاڼې په منځ
کښې «حالت (۲): $x \in C$»۔
وروسته د پاڼې پاسنۍ او لاندنۍ
نيمايي د ثبوت په ګامونو ډکې کړئ۔

د حالتونو له مخې ثبوت په ځانګړې توګه هغه
وخت ګټور دے چې غوښتل شوې پايله پخپله په
فصل کښې وي۔ دا يې يوه ساده بېلګه ده:

\begin{prop}
فرض کړئ $B \subseteq D$ او $C \subseteq E$۔
نو $B \cup C \subseteq
D \cup E$۔
\end{prop}

\begin{proof}
  فرض کړئ (الف) $B \subseteq D$ او (ب)
  $C \subseteq E$۔ د تعريف له مخې،
  هر $x \in B$، $ \in D$ هم دے
  (ج)، او هر $x \in C$، $ \in E$
  هم دے (د)۔
  د $B \cup C \subseteq D \cup E$
  د ښودلو دپاره بايد وښيو چې که
  $x \in B \cup C$ وي، نو
  $x \in D \cup E$ وي (د
  $\subseteq$ له تعريفه)۔
  $x \in B \cup C$ هغه وخت او
  يوازې هغه وخت وي چې $x \in B$
  يا $x \in C$ وي (د~$\cup$
  له تعريفه)۔ همدارنګه،
  $x \in D \cup E$ هغه وخت او
  يوازې هغه وخت وي چې $x \in D$
  يا $x \in E$ وي۔ نو بايد وښيو:
  د هر $x$ دپاره، که $x \in B$
  يا $x \in C$ وي، نو $x \in D$
  يا $x \in E$ وي۔

  \begin{quote}
  تر دې ځايه مو يوازې تعريفونه پرانيستي
  دي!{} خپل بيان مو د $\subseteq$ او
  $\cup$ له يادولو پرته بيا وايه، او اوس
  بايد يوه عمومي شرطي خبره ثابته کړو۔
  د پورته بحث له مخې، د دې دپاره
  $x$ داسې هر ډول څيز نيسو چې د
  «که» برخه يې رښتيا فرضوو؛ بيا
  ښيو چې د «نو» برخه هم رښتيا ده۔
  په بل عبارت، فرضوو چې $x
  \in B$ يا $x \in C$، او ښيو
  چې $x \in D$ يا $x \in
  E$۔\footnote{دا پراګراف يوازې
    زموږ کار تشريح کوي؛ د ثبوت
    برخه نه ده۔ د خپلو ثبوتونو
    د ليکلو پر مهال دومره تفصيل
    اړين نه دے۔}
  \end{quote}

  فرض کړئ $x \in B$ يا $x \in C$۔
  بايد وښيو چې $x \in D$ يا
  $x \in E$۔ حالتونه بېلوو۔

  حالت ۱: $x \in B$۔ د (ج)
  له مخې، $x \in D$۔ ځکه نو،
  $x \in D$ يا $x \in
  E$۔ (دلته مو په مخکېنۍ
  فرعي برخه کښې څېړل شوے
  استنتاج وکړ!)

  حالت ۲: $x \in C$۔ د (د)
  له مخې، $x \in E$۔ ځکه نو،
  $x \in D$ يا $x \in E$۔
 \end{proof}

\subsection{د وجود د دعوې ثابتول}

کله چې د وجود د دعوې د ثابتولو غوښتنه
درڅخه کېږي، پوښتنه عموماً داسې وي:
«ثابت کړئ چې داسې يو~$x$ شته چې
$\dots x \dots$»؛ يعنې داسې يو
څيز شته چې د «$\dots x
\dots$» په عبارت کښې راغلې
ځانګړنه لري۔ دلته بايد مناسب
څيز وټاکئ او وښيئ چې غوښتل
شوې ځانګړنه لري۔ دا کار
ساده ښکاري، خو دا ډول ثبوت
ستونزمن کېدے شي۔ عموماً بايد
يو څيز \emph{جوړ} يا
\emph{تعريف} کړئ او بيا
ثابته کړئ چې هغه ټاکلې
ځانګړنه لري۔ د سم څيز
موندل، د ځانګړنې ثابتول،
او کله کله ان دا ښودل
چې څيز مو په رښتيا تعريف
کړے دے، ټول ګران کېدے شي!{}

عموماً دا کار د څيز په ټاکلو ليکئ،
لکه «$x$ دې \dots وي» (دلته
\dots{} څرګندوي چې کوم څيز
مو په پام کښې دے)۔ کېدے شي
لومړے ثابته کړئ چې $\dots$
په رښتيا يو موجود څيز ښيي؛
بيا وښيئ چې $x$ د~$Q$
ځانګړنه لري۔ دا يې يوه ساده
بېلګه ده۔

\begin{prop}
  فرض کړئ $x \in B$۔ نو داسې يو~$A$
  شته چې $A \subseteq
  B$ او $A \neq \emptyset$۔
\end{prop}

\begin{proof}
  فرض کړئ $x \in B$۔ $A = \{x\}$ ونيسئ۔
  \begin{quote}
    دلته مو د~$A$ سټ د
    !!{element}s په شمېرلو تعريف کړے دے۔
    څرنګه چې $x$ يو څيز
    فرضوو، او تل د !!{element}s
    په شمېرلو سټ جوړولے
    شو، دلته دا جلا ثابتول
    نه دي پکار چې د~$A$
    سټ مو په برياليتوب
    تعريف کړے دے۔ خو دا
    لا بايد وښيو چې $A$
    د بيان غوښتل شوې
    ځانګړنې لري؛ بې له
    دې ثبوت بشپړ نه دے!{}
  \end{quote}
  څرنګه چې $x \in A$،
  نو $A \neq \emptyset$۔
  \begin{quote}
    دا د $A$ له تعريفه
    د $\{x\}$ په توګه
    او له ښکاره واقعيتونو
    څخه راځي چې
    $x \in \{x\}$
    او $x \notin \emptyset$۔
  \end{quote}
  ځکه چې $x$ د~$\{x\}$
  يوازينے !!{element} دے،
  او $x \in B$، د~$A$
  هر !!{element} د~$B$
  !!a{element} هم دے۔
  د~$\subseteq$ له تعريفه،
  $A \subseteq B$۔
\end{proof}

\subsection{د وجود د دعوو کارول}

فرض کړئ پوهېږئ چې يوه وجودي دعوه
رښتيا ده (يا مو ثابته کړې، يا يې
د کارولو فرض لرئ)، لکه «د ځينو~$x$
دپاره $x \in A$» يا «يو
$x \in A$ شته»۔ که يې په ثبوت
کښې کاروئ، داسې چلند وکړئ لکه
د هغو څيزونو له ډلې د يوه نوم
درسره وي چې فرض يې شتون ښيي۔
څرنګه چې په $A$ کښې لږ تر لږه
يو څيز شته، نوم پرې ايښودل
ممکن دي۔ ښايي تاسو ونه شئ
کړے هغه په ځانګړې توګه
وټاکئ يا يې، له دې پرته چې
$\in A$ دے، نور تشريح کړئ۔
خو د ثبوت دپاره يې د ټاکلو
په توګه ونيسئ او يو نوم
ورکړئ۔ مهمه ده چې داسې نوم
وټاکئ چې مخکې مو نه وي
کارولے او په فرضونو کښې
هم نه وي راغلے، ګنې تېروتنه
پېښېدے شي۔ په ثبوت کښې دا
ګام د «د ځينو $x$ دپاره
$x \in A$» څخه «$a
\in A$ دې وي» ته په تګ
ښيئ۔ اوس د~$a$ په اړه
استدلال او نور فرضونه
کارولے شئ، تر څو يوې
پايلې $p$ ته ورسېږئ۔
که $p$ نور~$a$ نه يادوي،
نو $p$ د $a \in A$
له فرضه خپلواکه ده؛ ښودلې
مو ده چې يوازې د «د ځينو
$x$ دپاره $x \in A$»
له فرضه راځي۔

\begin{prop}
که $A \neq \emptyset$ وي، نو
$A \cup B \neq \emptyset$۔
\end{prop}

\begin{proof}
  فرض کړئ $A \neq \emptyset$۔
  نو د ځينو $x$ دپاره $x \in A$۔
  \begin{quote}
    دلته مو لومړے د بيان
    فرض بيا ووايه۔ دا فرض،
    يعنې $A \neq \emptyset$،
    يوه وجودي دعوه پټوي
    چې د څو تعريفونو له
    پرانيستلو وروسته
    راښکاره کېږي۔ د $=$
    تعريف وايي چې $A = \emptyset$
    هغه وخت او يوازې هغه
    وخت وي چې هر $x \in A$
    په $\in \emptyset$ کښې
    هم وي او هر $x \in
    \emptyset$ په $\in A$
    کښې هم وي۔ د دواړو
    خواوو په نفي سره،
    $A \neq
    \emptyset$ هغه وخت او
    يوازې هغه وخت وي چې
    يا ځينې $x \in A$ په
    $\notin \emptyset$ کښې وي،
    يا ځينې $x \in \emptyset$
    په $\notin A$ کښې وي۔
    څرنګه چې هېڅ شے په
    $\in \emptyset$ کښې نه
    دے، دويم بديل هېڅکله
    رښتيا نه شي کېدے؛ او
    «$x \in A$ او $x
    \notin \emptyset$» يوازې
    $x \in A$ ته راټيټېږي۔
    نو $x \neq
    \emptyset$ هغه وخت او
    يوازې هغه وخت وي چې
    د ځينو $x$ دپاره $x \in A$۔
    \emph{د سرچينې يادونه
    (OLSTH-033): په دغه
    چاپي جمله کښې د نه
    تشوالي په کيڼ اړخ
    کښې د څيز نښه راغلې،
    خو د ثبوت دعوه د
    سټ د نه تشوالي ده؛ اصلي
    فورمول نه دے بدل شوے۔}
    دا وجودي دعوه ده۔
    اوس د~$A$ له
    !!{element}s څخه
    يوه ته نوم ورکوو:
  \end{quote}
  $a \in A$ دې وي۔
  \begin{quote}
    اوس مو په $\in A$
    کښې يوه څيز ته نوم
    ورکړ۔ د~$a$ په
    اړه استدلال ته دوام
    ورکوو، خو پام کوو چې
    له $a \in A$ پرته
    نور څه فرض نه کړو:
  \end{quote}
  ځکه چې $a \in A$،
  $a \in A \cup B$،
  د~$\cup$ له تعريفه۔
  نو د ځينو $x$ دپاره
  $x \in A \cup B$،
  يعنې $A \cup B \neq \emptyset$۔
  \begin{quote}
    په دغه وروستي ګام
    کښې له «$a \in A \cup B$»
    څخه «د ځينو $x$
    دپاره $x \in A \cup B$»
    ته لاړو۔ وروستۍ
    خبره نور $a$ نه يادوي؛
    نو پوهېږو چې «د ځينو
    $x$ دپاره $x \in A \cup B$»
    يوازې له «د ځينو
    $x$ دپاره $x \in A$»
    څخه راځي۔ همدا معنا
    لري چې $A \cup B \neq
    \emptyset$۔
  \end{quote}
\end{proof}

ښايي ښه تمرين وي چې تړلي متغيرونه،
لکه «$x$»، له فرضي نومونو،
لکه $a$، بېل وساتو؛ موږ هم
همداسې وکړل۔ خو په عمل کښې
ډېر ځله دا بېلوالے نه ساتو
او $x$ کاروو، لکه دلته:
\begin{quote}
فرض کړئ $A \neq \emptyset$،
يعنې يو $x \in A$ شته۔
د $\cup$ له تعريفه،
$x \in A \cup B$۔
نو $A \cup B \neq \emptyset$۔
\end{quote}
خو په دې توګه استدلال
کښې بايد په ډېر پام
د بېلابېلو وجودي دعوو
دپاره بېل $x$ او $y$
وکاروئ۔ د بېلګې په توګه،
لاندې د «که $A \neq
\emptyset$ او $B \neq \emptyset$
وي، نو $A \cap B \neq \emptyset$»
ثبوت \emph{سم نه دے} (ځکه
دا دعوه رښتيا نه ده):
\begin{quote}
فرض کړئ $A \neq \emptyset$
او $B \neq \emptyset$۔
نو د ځينو $x$ دپاره
$x \in A$، او هم د ځينو
$x$ دپاره $x \in B$۔
څرنګه چې $x \in A$
او $x \in
B$، $x \in A \cap B$،
د~$\cap$ له تعريفه۔
نو $A \cap B \neq
\emptyset$۔
\end{quote}
ايا هغه ناسم ګام پېژندلے
شئ او ويلے شئ چې دا
پايله ولې نه راځي؟

\end{document}
